\bookchapter{Git and GitHub}\label{ch:git3}

\section*{Explanation}
So far all history lived on your computer. To work in a team, to have a backup and to get the code onto the robot, you need a second copy on a server: a \textbf{remote}. GitHub is the best known place to host remotes. This chapter shows the commands that move commits between your repository and a remote, and the working style around GitHub: branches, pull requests, reviews and automatic tests.

\section*{Remotes}
A \textbf{remote} is a named link to another copy of the repository. By convention the main one is called \texttt{origin}.
\begin{center}
\begin{tikzpicture}[box/.style={draw,rounded corners,thick,minimum width=2.9cm,minimum height=1.5cm,align=center},>={Stealth[length=2.5mm]},lab/.style={font=\scriptsize\ttfamily,fill=white,inner sep=1.5pt}]
  \node[box,fill=blue!10] (w) at (0,0) {your\\working directory};
  \node[box,fill=green!12] (l) at (5.2,0) {your local\\repository\\{\footnotesize branches: \texttt{main}}\\{\footnotesize \texttt{origin/main}}};
  \node[box,fill=orange!15] (g) at (10.4,0) {\textbf{GitHub}\\\texttt{origin}\\{\footnotesize branch \texttt{main}}};
  \draw[->,thick] (w.20) -- node[lab,above]{commit} (l.160);
  \draw[->,thick] (l.20) -- node[lab,above]{push} (g.160);
  \draw[->,thick] (g.200) -- node[lab,below]{fetch} (l.340);
  \draw[->,thick,dashed] (l.200) -- node[lab,below]{merge/rebase} (w.340);
  \node[font=\footnotesize,gray!60!black,align=center] at (7.8,1.7) {\texttt{git pull} = fetch + merge (or rebase)};
\end{tikzpicture}
\captionof{figure}{Where the commands act. \texttt{origin/main} in your local repository is your \emph{memory} of the remote branch as it was at the last fetch or push.}\label{fig:g3_remote}
\end{center}

\subsection*{Getting a repository}
\begin{verbatim}
$ git clone git@github.com:ada/walker.git      # SSH URL
$ git clone https://github.com/ada/walker.git  # HTTPS URL
$ cd walker
$ git remote -v
origin	git@github.com:ada/walker.git (fetch)
origin	git@github.com:ada/walker.git (push)
\end{verbatim}
\texttt{git clone} copies the whole history and sets up \texttt{origin}. For a repository that you started yourself with \texttt{git init}, create an empty repository on GitHub (button \emph{New}), and connect it:
\begin{verbatim}
$ git remote add origin git@github.com:ada/walker.git
$ git push -u origin main
\end{verbatim}
The \texttt{-u} remembers that your \texttt{main} follows \texttt{origin/main}, so later a plain \texttt{git push} and \texttt{git pull} are enough.

\subsection*{The four commands}
\begin{description}[nosep]
  \item[\texttt{git fetch}] downloads new commits and updates \texttt{origin/main}. It never touches your files; it is always safe.
  \item[\texttt{git pull}] is \texttt{fetch} followed by merging the remote branch into your current branch.
  \item[\texttt{git push}] uploads your new commits.
  \item[\texttt{git status -sb}] tells you how far you are ahead or behind.
\end{description}
\begin{verbatim}
$ git fetch
From github.com:ada/walker
   088b1b2..1f879d0  main       -> origin/main
$ git status -sb
## main...origin/main [ahead 1, behind 1]
\end{verbatim}
``Ahead 1, behind 1'' means: you have one commit that GitHub does not have, and GitHub has one that you do not have; the histories have diverged.

\subsection*{The rejected push}
\begin{verbatim}
$ git push
 ! [rejected]        main -> main (fetch first)
error: failed to push some refs to 'github.com:ada/walker.git'
hint: Updates were rejected because the remote contains work that you do
hint: have locally. ...
\end{verbatim}
This is the most common Git message in a team. A colleague pushed first. Nothing is broken: integrate their work, then push again.
\begin{verbatim}
$ git pull --rebase        # replay my commits on top of theirs
Successfully rebased and updated refs/heads/main.
$ git push
\end{verbatim}
A plain \texttt{git pull} would work too, but creates a merge commit ``Merge branch main of github.com...'' every time, which clutters the log. Make rebase the default: \verb!git config --global pull.rebase true!. Never use \texttt{git push --force} to get around the rejection: it destroys the work of the colleague (Chapter~\ref{ch:git4}).

\section*{Authentication}
GitHub must know who pushes. Two ways work well.
\begin{description}
  \item[SSH keys (recommended).] You already know them from Chapter~\ref{ch:linux3}. Create a key with \verb!ssh-keygen -t ed25519!, copy the \emph{public} key (\verb!cat ~/.ssh/id_ed25519.pub!), and add it in GitHub under \emph{Settings $\rightarrow$ SSH and GPG keys $\rightarrow$ New SSH key}. Test it:
\begin{verbatim}
$ ssh -T git@github.com
Hi ada! You've successfully authenticated, but GitHub does not
provide shell access.
\end{verbatim}
  Use the SSH form of the URL. Change an existing clone with \verb!git remote set-url origin git@github.com:ada/walker.git!.
  \item[HTTPS with a token.] A password is not accepted for Git any more; you need a \emph{personal access token} or the GitHub command line tool \texttt{gh} (\verb!sudo apt install gh!, then \verb!gh auth login!) which stores the login for you.
\end{description}
\textbf{On the robot} use a \emph{deploy key}: create a key pair on the robot (without personal rights), and add the public key to the repository under \emph{Settings $\rightarrow$ Deploy keys} (read-only). Then the robot can \texttt{git pull} the code, but cannot change it, and it does not carry your personal identity. Never copy your personal private key to a robot that many people use.

\section*{The GitHub flow}
A simple, widely used working style for teams (Figure~\ref{fig:g3_flow}):
\begin{center}
\begin{tikzpicture}[box/.style={draw,rounded corners,thick,minimum width=2.5cm,minimum height=1.1cm,align=center,font=\footnotesize},>={Stealth[length=2.5mm]}]
  \node[box,fill=blue!10] (a) at (0,0) {1. branch\\from \texttt{main}};
  \node[box,fill=blue!10] (b) at (3.1,0) {2. commit\\and push};
  \node[box,fill=orange!15] (c) at (6.2,0) {3. open a\\pull request};
  \node[box,fill=orange!15] (d) at (9.3,0) {4. review\\and tests};
  \node[box,fill=green!12] (e) at (12.4,0) {5. merge\\into \texttt{main}};
  \draw[->,thick] (a)--(b); \draw[->,thick] (b)--(c); \draw[->,thick] (c)--(d); \draw[->,thick] (d)--(e);
  \draw[->,thick,red!70!black] (d.south) -- ++(0,-0.7) -| node[pos=0.25,below,font=\scriptsize]{changes requested: commit and push again} (b.south);
  \draw[->,thick,gray!60!black] (e.north) -- ++(0,0.8) -| node[pos=0.25,above,font=\scriptsize]{pull \texttt{main}, delete the branch, start again} (a.north);
\end{tikzpicture}
\captionof{figure}{The GitHub flow. \texttt{main} always contains working code; all changes arrive through short-lived branches and pull requests.}\label{fig:g3_flow}
\end{center}
\begin{verbatim}
$ git switch main && git pull               # 1. start up to date
$ git switch -c fix/camera-timeout
$ # edit, test, commit (small commits)      # 2.
$ git push -u origin fix/camera-timeout
remote: Create a pull request for 'fix/camera-timeout' on GitHub by
remote: visiting: https://github.com/ada/walker/pull/new/fix/camera-timeout
\end{verbatim}
GitHub prints a link; open it to create the \textbf{pull request} (PR). A PR is a proposal: ``please merge my branch into \texttt{main}''. It is the place where the team discusses the change.

\subsection*{A good pull request}
\begin{itemize}[nosep]
  \item A \textbf{title} that says what changes (``Fix camera timeout after reconnect'').
  \item A \textbf{description}: what and \emph{why}, how you tested, a link to the issue (\texttt{Closes \#12} closes the issue automatically on merge), a screenshot or log if useful.
  \item \textbf{Small}: a PR that one person can review in $15$ minutes. Large PRs get skimmed and bugs slip through. Split big work.
  \item Open it as \emph{draft} when you want early feedback on unfinished work.
\end{itemize}
\textbf{Reviewers} read the \emph{Files changed} tab, comment on single lines, and finally choose \emph{Approve} or \emph{Request changes}. Review the code, not the person: ask questions (``what happens if the camera returns \texttt{None}?''). As author, answer every comment, push fixes as new commits; the PR updates itself.

\subsection*{Merging a pull request}
GitHub offers three buttons:
\begin{center}
\renewcommand{\arraystretch}{1.25}
\begin{tabular}{p{3.4cm}p{9.8cm}}
\textbf{Option} & \textbf{Result}\\ \hline
Create a merge commit & all branch commits stay, plus a merge commit. Full history, but noisy.\\
Squash and merge & all commits of the branch become \emph{one} commit on \texttt{main}. Clean history; the messy ``fix typo'' commits vanish. A good default for small PRs.\\
Rebase and merge & commits are replayed onto \texttt{main}, linear, no merge commit.\\
\end{tabular}
\end{center}
After the merge, delete the branch (GitHub offers a button) and update your own copy:
\begin{verbatim}
$ git switch main
$ git pull
$ git branch -d fix/camera-timeout
$ git fetch --prune          # forget remote branches deleted on GitHub
\end{verbatim}

\section*{Protecting main}
In the repository settings (\emph{Settings $\rightarrow$ Branches}, a \emph{branch protection rule} or \emph{ruleset}) you can enforce the flow: nobody pushes directly to \texttt{main}, a PR needs at least one approval, and the automatic tests must pass. This is cheap insurance for the code that runs on a robot.

\section*{Automatic tests: GitHub Actions}
\textbf{GitHub Actions} runs a program on GitHub's servers for every push or pull request, for example your tests. You describe it in a YAML file in the folder \texttt{.github/workflows/}:
\begin{verbatim}
# .github/workflows/ci.yml
name: CI
on:
  push:
    branches: [main]
  pull_request:
jobs:
  test:
    runs-on: ubuntu-latest
    steps:
      - uses: actions/checkout@v4
      - uses: actions/setup-python@v5
        with:
          python-version: "3.12"
      - run: pip install -r requirements.txt
      - run: python -m pytest
\end{verbatim}
The result appears on the PR as a green check or a red cross, and the merge button can be blocked while it is red. Typical jobs: run tests, check the code style (\texttt{ruff}), build a package. Indentation is part of YAML: use spaces, never tabs. (Use a pinned action version like \texttt{@v4} and check the current one in the action's documentation.)

\section*{Issues, forks and the files around}
\begin{description}
  \item[Issues] are the to-do list and bug tracker of a repository. One issue per problem, with steps to reproduce and the expected behaviour. Labels (\texttt{bug}, \texttt{enhancement}) and assignees organise them.
  \item[Fork.] If you may not push to a repository (open-source project), \emph{fork} it: GitHub copies it into your account. You push to your fork and open a PR from it. Keep it current with a second remote:
\begin{verbatim}
$ git remote add upstream git@github.com:team/walker.git
$ git fetch upstream
$ git merge upstream/main
\end{verbatim}
  \item[Files that every repository should have:] \texttt{README.md} (what it is, how to install and run, in five minutes), \texttt{LICENSE} (what others may do with your code), \texttt{.gitignore}, and optionally \texttt{CONTRIBUTING.md} (how to propose changes) and \texttt{.github/CODEOWNERS} (who must review which folder).
  \item[Releases.] Tag a version (\verb!git tag -a v1.0 -m "..."!, \verb!git push origin v1.0!) and publish it under \emph{Releases}. On the robot you then check out a \emph{known} version instead of whatever was on \texttt{main} yesterday.
\end{description}

\section*{Secrets and safety}
\begin{keyidea}[A pushed secret is a leaked secret]
Never commit passwords, API tokens or private keys. Bots scan GitHub within minutes. If it happens, \textbf{revoke the secret immediately} (change the password, delete the token); removing the file from history is only the second step. Keep secrets in environment variables, in an \texttt{.env} file listed in \texttt{.gitignore}, or in the \emph{Secrets} feature of GitHub Actions (\emph{Settings $\rightarrow$ Secrets and variables}).
\end{keyidea}
Also remember that a public repository is public for the whole world, and a ``private'' repository is private only as long as the access list is right.

\begin{example}
Ada fixes a bug: \texttt{git switch -c fix/camera-timeout}, two commits, \texttt{git push -u origin fix/camera-timeout}, PR with ``Closes \#12''. The CI check is green, Bob asks for a log message, Ada pushes one more commit. Bob approves, Ada uses \emph{Squash and merge}, deletes the branch on GitHub, and runs \texttt{git switch main \&\& git pull \&\& git branch -d fix/camera-timeout}. On the robot, a one-line \texttt{ssh k1 'cd walker \&\& git pull --ff-only'} brings in the fix.
\end{example}

\section*{Practice problems}
\begin{enumerate}[label=\textbf{\thechapter.\arabic*}]
  \item What is the difference between \texttt{git fetch} and \texttt{git pull}? Which one can never change your working files?
  \item You created a local repository and want to publish it to the empty GitHub repository \texttt{ada/walker}. Write the commands (SSH URL).
  \item \texttt{git status -sb} shows \verb!## main...origin/main [ahead 2]!. What does it mean and what do you do?
  \item Your push is rejected with \texttt{fetch first}. Why, and what is the clean way to fix it?
  \item Describe the five steps of the GitHub flow.
  \item List the steps to let your laptop authenticate to GitHub with SSH.
  \item Why is a deploy key a better choice for the robot than your personal key?
  \item What is the difference between \emph{Squash and merge} and \emph{Create a merge commit}? When is squash a good choice?
  \item Write a PR description for a change that increases a timeout from $2$ to $5$ seconds to fix issue 12.
  \item[\textbf{\thechapter.10}$\star$] You accidentally pushed a file with an API token to a public repository. List what you do, in the right order.
\end{enumerate}
